\scanpage{155}
\mspage{13}
i.e. $X_{/[x]}\cap X_{/[y]}\ne\varnothing$ [où $[x]$ est \deleted{\unsure{l'id$(x)$}} \inserted{$x$} considéré comme \deleted{objet} \inserted{élément} de l'ens ordonné $\mathrm{Ob}\,X$].
Il \unsure{revient aussi|revient encore} au $=$ de dire que les catégories induites $\deleted{\unsure{I}}X_{/x}$ sont \unsure{irréductible|irréductibles}.
% slanted
\margin{faux cf \unsure{aussi le|aussi} \unsure{contre-exemple} \unsure{précédent|précédente}}

\deleted{\unsure{Comme}} \uline{Exemple} Une catégorie \uline{ayant un objet initial}, ou seulement un ``plus petit objet'', est irréductible.
\deleted{\unsure{Donc aussi|D'une manière} toute catégorie dont tte composante connexe admet un plus petit objet, est loc. irréductible.}
Inversement, si l'ens \uline{ordonné} $\mathrm{Ord}(X)$ associé à $X$ [quotient de $\mathrm{Ob}\,X$ par la relation d'équivalence associée à son préordre] est fini, p. ex. $X$ fini, \deleted{donc} ou seulement loc fini, p. ex. $X$ loc. fini (i.e. les catégories $X_{/x}$ finies), alors si $X$ est irréductible, $\mathrm{Ob}\,X$ a un plus petit élément i.e. $X$ a un plus petit objet.
C'est vrai plus gén. encore, si $\mathrm{Ord}(X)$ est artinien, i.e. \deleted{toute} toute suite décroissante dans $\mathrm{ord}(X)$ est stationnaire. //
Mais la condition intéressante est la condition \unsure{toposique} de \uline{locale} irréductibilité, qui équivaut à $I$ pseudo-filtrante \inserted{\unsure{à g.}}, i.e. les comp. connexes de $I$ \deleted{pseudo-}filtrantes, i.e. les conditions $\mathrm{PS}1$, $\mathrm{PS}2$
\uline{PS1} Tout diagramme
\diagram{\begin{tikzcd}[ampersand replacement=\&,row sep=0pt,column sep=small]
 x\arrow[dr] \& {}\\ {} \& s\\ y\arrow[ur] \& {}
\end{tikzcd}}
se complète en
\diagram{\begin{tikzcd}[ampersand replacement=\&,row sep=0pt,column sep=small]
 \& x\arrow[dr] \& {}\\ z\arrow[ur]\arrow[dr] \& {} \& s\\ \& y\arrow[ur] \& {}
\end{tikzcd}}
\uline{commutatif}

\uline{PS2} Pour \unsure{tte|tt} double flèche $x\rightrightarrows y$, existe un égalisateur
\diagram{\begin{tikzcd}[ampersand replacement=\&,column sep=small]
 z\arrow[r] \& x\arrow[r,shift left]\arrow[r,shift right] \& y
\end{tikzcd}}
\deleted{\unsure{est} \unsure{commutatif} \unsure{ou} $\unsure{z\to x}$ \unsure{(égalise)} \illegible{1}}
% archive sheet number: 155
